
In this paper, we study the impact on interpretability of adapting a Large Pre-Trained Model $L$ using an adapter $A$ using a mechanistic interpretability tool $S$. Practically, we restrict attention to visual encoders CLIP, SigLIP and ALIGN that are part of LVMs and DINO which is a pure vision encoder. We restrict this study to Sparse Auto-Encoders (SAEs) for interpretability and to LoRA, DoRA, MaPLe, and TAP as adapters. We believe our conclusions will hold for other combinations of encoders, adapters and interpretability tools.

An LPTM $L$ is adapted to a new dataset using an adapter $A$ to yield $\hat{L}$. $S_L$ is an SAE trained on activations of a specific layer of $L$. $S_L(x)$ is its decoder output for input $x$ and $S_L^E(x)$ is its {\em encoding} at the bottleneck layer. An image $I$ is transformed to $L(I)$ by the model and maps to $S_L(L(I))$ by the SAE.

SAEs trained using large, pretraining-like corpus can capture a generic dictionary. They could be rich enough to explain an adapted $\hat L$. We call this a {\em Generic SAE}, denoted by $\bar S_L$. $S_{\hat L}$ is an SAE trained on the adapted model. We call this a {\em Specific SAE}, denoted by $S_{\hat L}$. We explore the explainability of $\bar S_L$ and $S_{\hat L}$ using several standard metrics to understand the behavior of Generic SAE and Specific SAEs as the distribution moves farther from the corpus.

We study different metrics to analyze interpretability while using adapters. $S_L$ and $S_{\hat{L}}$ trained before and after adaptation. Specifically, we ask the question: can concepts learned by $S_L$ explain how the adapted model $\hat L$ behaves? We study different metrics on $S_L(L(I)$), $S_L(\hat{L}(I))$ and $S_{\hat L}(\hat{L}(I))$.

\begin{table}[h]
    \centering
    \begin{tabular}{|c||c|c|c|c|c|}
    \hline
        LPTM & CLIP & SigLIP & DINO & ALIGN & \\
    \hline
        Adapter & LoRA & DoRA & MaPLe & TAP & \\
    \hline
        SAEs & Gated & Top-K & Matryoshka && \\
    \hline
    \end{tabular}
    \caption{Configurations explored}
    \label{configTable}
\end{table}

We study the behavior under different combinations of LPTMs, adapters and SAEs as given in Table \ref{configTable}. Experiments are carried out on datasets commonly used for doman adaptation and domain generalization literature. The datasets are given in Table \ref{datasetTable}.


\begin{table}[h]
    \centering
    \begin{tabular}{|c|c|c|}
    \hline
       NaturalShift  &  & \\
    \hline
       Medical  &  & \\
    \hline
       Others  &  & \\
    \hline
    \end{tabular}
    \caption{Caption}
    \label{datasetTable}
\end{table}


\subsection{Tool: sparse autoencoders across architectures}
\label{sec:method-tool}

Rather than committing to a single SAE architecture, we sweep across
several variants---ReLU, Top-$k$, Gated, Matryoshka, and PatchSAE---to
test whether our conclusions about SAE behaviour under adaptation hold
across dictionary-learning design choices rather than being artefacts of
one parameterisation. Each variant trades off sparsity control, stability,
and feature granularity differently, and a finding that survives this
sweep is a finding about adaptation itself rather than about a particular
objective. We foreground PatchSAE in the exposition because it operates
on patch tokens and therefore supports both local analysis---textures,
artefacts, acquisition noise that dominate many domain shifts---and their
global composition into semantic categories; the remaining variants are
trained and evaluated under the same protocol and reported alongside.

\paragraph{Common skeleton.}
All variants share a linear encoder, a nonlinearity $\sigma(\cdot)$, and
a linear decoder trained to reconstruct hooked ViT token activations. Let
$z \in \mathbb{R}^d$ be a hooked token vector (CLS or patch) at layer
$\ell$, and let $d_{\mathrm{SAE}} \gg d$ be the overcomplete latent
dimension. The general form is
\begin{equation}
\mathrm{SAE}(z) \;=\; W_D^{\top}\,h(z) + b_D,
\qquad
h(z) \;=\; \sigma\!\left(W_E^{\top}(z - b_D) + b_E\right)
\in \mathbb{R}^{d_{\mathrm{SAE}}}_{\geq 0},
\label{eq:sae-general}
\end{equation}
where $W_E \in \mathbb{R}^{d \times d_{\mathrm{SAE}}}$,
$W_D \in \mathbb{R}^{d_{\mathrm{SAE}} \times d}$, and $b_E, b_D$ are
bias terms. The variants differ in the choice of $\sigma$ and in the
sparsity-inducing term added to the reconstruction loss.

\paragraph{ReLU SAE~\citep{cunningham2024sparse}.}
The simplest variant uses $\sigma = \mathrm{ReLU}$ with an $\ell_1$
penalty on the latent activations:
\begin{equation}
\mathcal{L}_{\mathrm{ReLU}}
\;=\;
\|\mathrm{SAE}(z) - z\|_2^2
\;+\;
\lambda_{\ell_1}\,\|h(z)\|_1.
\label{eq:relu-sae}
\end{equation}
The nonnegative parameterisation is easy to optimise but can leave a
fraction of latents dead or underutilised.

\paragraph{Top-$k$ SAE~\citep{gao2025scaling}.}
Top-$k$ SAEs impose exact sparsity by retaining only the $k$ largest
pre-activations and zeroing the rest:
\begin{equation}
h(z) \;=\; \mathrm{TopK}_{k}\!\left(W_E^{\top}(z - b_D) + b_E\right),
\label{eq:topk-sae}
\end{equation}
where $\mathrm{TopK}_{k}(\cdot)$ keeps the $k$ entries with the largest
values and sets the remainder to zero. No explicit sparsity term is
needed, so the loss reduces to reconstruction alone:
\begin{equation}
\mathcal{L}_{\mathrm{TopK}} \;=\; \|\mathrm{SAE}(z) - z\|_2^2.
\label{eq:topk-loss}
\end{equation}

\paragraph{Gated SAE~\citep{rajamanoharan2024improving}.}
Gated SAEs decouple the decision \emph{whether} a latent fires from the
decision \emph{how strongly} it fires, using separate gating and
magnitude pathways that share a tied encoder direction. Writing
$\pi(z) = W_E^{\top}(z - b_D)$, the gated activation is
\begin{equation}
h_{\mathrm{gate}}(z)
\;=\;
\mathbb{1}\!\left(\pi(z) + b_{\mathrm{gate}} > 0\right)
\odot
\mathrm{ReLU}\!\left(\exp(r)\odot\pi(z) + b_{\mathrm{mag}}\right),
\label{eq:gated-sae}
\end{equation}
where $r \in \mathbb{R}^{d_{\mathrm{SAE}}}$ is a learned rescaling and
$b_{\mathrm{gate}}, b_{\mathrm{mag}}$ are bias terms. The loss combines
reconstruction, an $\ell_1$ penalty on the gate pre-activation, and an
auxiliary reconstruction term that prevents the gate path from collapsing:
\begin{equation}
\mathcal{L}_{\mathrm{Gated}}
\;=\;
\|\mathrm{SAE}(z) - z\|_2^2
\;+\;
\lambda_{\ell_1}\,\|\mathrm{ReLU}(\pi(z) + b_{\mathrm{gate}})\|_1
\;+\;
\|\hat{z}_{\mathrm{gate}}(z) - z\|_2^2,
\label{eq:gated-loss}
\end{equation}
where $\hat{z}_{\mathrm{gate}}$ is the reconstruction produced using only
the gate pathway. This decoupling improves stability and reduces
feature-shrinkage bias relative to ReLU SAEs.

\paragraph{Matryoshka SAE~\citep{bussmann2025learning,zaremba2025interpreting}.}
Matryoshka SAEs enforce a nested hierarchy on the dictionary by training
a sequence of prefix-truncated reconstructions. Given an ordered set of
prefix sizes $m_1 < m_2 < \cdots < m_M = d_{\mathrm{SAE}}$, let
$h_{1:m}(z) \in \mathbb{R}^{d_{\mathrm{SAE}}}$ denote the latent vector
with only its first $m$ entries retained and the rest zeroed. The loss
sums reconstruction and sparsity terms over all prefixes:
\begin{equation}
\mathcal{L}_{\mathrm{Mat}}
\;=\;
\sum_{k=1}^{M} \alpha_k
\left[
\left\|W_D^{\top} h_{1:m_k}(z) + b_D - z\right\|_2^2
+ \lambda_{\ell_1}\,\|h_{1:m_k}(z)\|_1
\right],
\label{eq:matryoshka-loss}
\end{equation}
with weights $\alpha_k > 0$. The earliest latents are pressured to carry
coarse, broadly-used structure and later latents to refine it, yielding a
multi-granularity dictionary.

\paragraph{PatchSAE~\citep{lim2024sparse}.}
PatchSAE uses the ReLU objective of Eq.~\eqref{eq:relu-sae} but is
trained on \emph{patch tokens} rather than only the CLS token, so each
latent is a detector defined over the full $N \times d$ patch grid. This
matters for our setting because domain shifts often manifest locally
(texture, artefact, acquisition noise) and only aggregate into a
semantic category through composition across patches. For image-level
analysis under the PatchSAE variant, we threshold token-wise activations:
for neuron $s$ on image $i$, patch-level and image-level activation
counts are
\begin{equation}
a_{ij}[s] \;=\; \mathbb{1}\!\left(h_{ij}[s] > \tau\right),
\qquad
a_i[s] \;=\; \sum_{j=1}^{N} a_{ij}[s],
\label{eq:activation-counts}
\end{equation}
for a small threshold $\tau > 0$. These counts define top-activating
images per neuron and support the dataset- and class-level summaries used
in the diagnostics of Section~\ref{sec:method-tests}. The other variants
are applied to CLS-token activations under their respective losses and
share the rest of the training protocol.

\subsection{Procedure: paired SAE training and cross-evaluation}
\label{sec:method-procedure}

The standard PatchSAE protocol trains a single SAE on ImageNet activations
from the base model and reuses it across downstream datasets---treating
cross-domain reuse as a silent assumption. Our procedure turns that
assumption into a testable claim. For every (encoder, adapter, dataset)
triple, we train two SAEs under identical hyperparameters:
\begin{enumerate}
    \item a \emph{base SAE} $S_L$ on activations $H^{\ell}(x)$ from the
          frozen encoder $L$, and
    \item an \emph{adapted SAE} $S_{\hat{L}}$ on activations
          $\hat{H}^{\ell}(x)$ from the adapted encoder $\hat{L}$.
\end{enumerate}
We then evaluate three decomposition pipelines on each dataset:
$S_L(L(I))$ (base SAE on base model---the reference),
$S_L(\hat{L}(I))$ (base SAE on adapted model---the reuse assumption), and
$S_{\hat{L}}(\hat{L}(I))$ (adapted SAE on adapted model---the oracle). The
gap between the second and third pipelines quantifies \emph{how much the
SAE itself must adapt} for the dictionary to remain faithful; the gap
between the first and second quantifies \emph{how far adaptation pushes
activations off the base dictionary}.

\subsection{Discovery tests: diagnostics that reveal SAE degradation}
\label{sec:method-tests}

To characterise the gap along multiple axes, we compute four representation
checks on each pipeline, each grounded in prior SAE literature rather than
introduced ad hoc.

\paragraph{(1) Predictive utility.}
Can class-discriminative information still be read out from the SAE code?
Following PatchSAE, we report classification accuracy obtained by replacing
the encoder's intermediate activations with their SAE reconstruction (or
equivalently, by using aggregated image-level SAE features as input to a
lightweight classifier). A drop relative to the oracle pipeline indicates
the dictionary no longer carries the information the adapted model uses.

\paragraph{(2) Reconstruction fidelity.}
The mean squared reconstruction error
\begin{equation}
E_{\mathrm{rec}} \;=\; \mathbb{E}\!\left[\|\mathrm{SAE}(z) - z\|_2^2\right]
\label{eq:recon-error}
\end{equation}
measured on held-out activations directly tests whether the basis spans the
representation. A high $E_{\mathrm{rec}}$ for $S_L(\hat{L}(I))$ with a low
$E_{\mathrm{rec}}$ for $S_{\hat{L}}(\hat{L}(I))$ is a direct signature of
dictionary drift.

\paragraph{(3) Sparsity and usage.}
For each latent $s$, we track two quantities from SAE: the activated
frequency
\begin{equation}
\mathrm{Freq}[s]
\;=\;
\frac{\left|\{i \in \mathcal{I} : h_i[s] > 0\}\right|}{|\mathcal{I}|},
\label{eq:freq}
\end{equation}
and the mean activation among activated samples,
\begin{equation}
\mathrm{MeanAct}[s]
\;=\;
\frac{1}{\left|\{i \in \mathcal{I} : h_i[s] > 0\}\right|}
\sum_{\substack{i \in \mathcal{I} \\ h_i[s] > 0}} h_i[s].
\label{eq:meanact}
\end{equation}
These give an auditable usage profile: adaptation may induce dead features,
diffuse activations, or concentration on a narrower set of latents.

\paragraph{(4) Selectivity.}
We measure how cleanly top-activating images for a latent align with a
single concept using two complementary scores.

\emph{Label entropy.} For latent $s$ with reference set $R_s$ (top-$k$
images by image-level activation), let
$\mathrm{sum}_c = \sum_{i \in R_s : y_i = c} h_i[s]$ and
$p_c = \mathrm{sum}_c / \sum_{c' \in \mathcal{C}} \mathrm{sum}_{c'}$. Then
\begin{equation}
\mathrm{Ent}[s]
\;=\;
-\sum_{c \in \mathcal{C}} p_c \log p_c.
\label{eq:entropy}
\end{equation}
Entropy near zero indicates the latent concentrates on a single label.

\emph{Monosemanticity.} Let $x_i$ be the CLIP image embedding of image $i$
and $z_i^{h} \geq 0$ the activation of latent $h$ on $i$. Weighting each
image pair by $w_{ij}^{h} = z_i^{h} z_j^{h}$,
\begin{equation}
\mathrm{MS}_h
\;=\;
\frac{\sum_{i<j} w_{ij}^{h} \cos(x_i, x_j)}{\sum_{i<j} w_{ij}^{h}}.
\label{eq:monosemanticity}
\end{equation}
Pairs where the latent fires strongly on both images dominate the weighted
average; pairs where it is inactive on at least one contribute zero. High
$\mathrm{MS}_h$ indicates the latent fires on a semantically coherent set.

We summarise each metric over the top-$K$ latents (we use $K = 10$) to
match reporting in Section~\ref{sec:experiments}.

\paragraph{Putting the tests together.}
The four checks probe complementary failure modes: classification tests
whether the basis is still \emph{usable}, reconstruction tests whether it
is still \emph{sufficient}, sparsity tests whether it is still
\emph{structured}, and entropy and monosemanticity test whether it is
still \emph{selective}. A base-SAE dictionary that survives adaptation
should score comparably on $S_L(\hat{L}(I))$ and $S_{\hat{L}}(\hat{L}(I))$
across all four; a dictionary that breaks down under adaptation will
diverge on at least one, and the pattern of divergence tells us how.